% =============================================================================
% Certus -- Section: Experimental setup (DRAFT, Week 8)
% =============================================================================

\section{Experimental Setup}
\label{sec:setup}

\subsection{Data}

\paragraph{SDNET2018.}
SDNET2018 \citep{dorafshan2018sdnet} contains 56{,}092 human-labelled
$256\times256$ images of concrete bridge decks, pavements and walls, of
which 8{,}484 (15\%) are cracked. Every filename encodes a physical
structure identifier, giving 230 real structures. It is our primary
dataset: its labels are human, and its majority of clean images makes the
auto-clear decision meaningful. We draw a fresh deterministic 70/15/15
stratified split (seed 0). Unless stated otherwise, experiments use a fixed
3000-image sample of the calibration split (457 cracked) as their pool.

\paragraph{CBRI.}
The CBRI building-defect dataset (Roboflow) contains 1{,}518 images of
whole building facades, labelled for crack, vegetation, corrosion and
dampness by GPT-4V bootstrapping. We hand-labelled all 175 images outside
its training split, blind to the bootstrap annotations. 172 contain a
defect. We use CBRI for the label-noise study (\S\ref{sec:label-noise})
and as an unseen target site for the shift study. A domain classifier on
frozen CLIP features separates CBRI from SDNET2018 with AUC $0.9996$.

\paragraph{Provenance.}
Both datasets are version-controlled with DVC. All splits are frozen in a
manifest whose content hash
(\texttt{aeaa352b\ldots}) is logged with every result, along with each
score's model version, the pool size, the number of positives and the seeds.

\subsection{Tuning set}

Fitted scores learn from a fixed tuning set: 500 cracked and 500 uncracked
SDNET2018 \emph{training} images. It is disjoint from every calibration
pool. We also report each score's AUC on it as a sanity check of the
score's signal.

\subsection{Protocol}

Every number is a mean over 100 seeds. Each seed draws a fresh random
partition of the pool into calibration and test rows; the seed is derived
from $(\text{base seed}, i)$, so any single seed can be replayed in
isolation. We report mean $\pm$ SD and, for coverage, a 95\% percentile
bootstrap interval over seeds.

\begin{itemize}
  \item \textbf{E1, validity.} Split CP (LAC, APS, RAPS) and baselines
  B1--B3, $\alpha \in \{0.01, 0.05, 0.10, 0.20\}$, $n_{\text{cal}} \in \{50,
  200, 1000\}$, scores S1 and S2, on the SDNET2018 pool and on all of CBRI.
  The gate requires mean coverage within $\pm 2$ points of $1-\alpha$.
  \item \textbf{Risk control.} CRC and LTT ($\delta = 0.05$),
  $\varepsilon \in \{0.10, 0.05, 0.02, 0.01\}$, joint and conditional,
  $n_{\text{cal}} = 2000$, which leaves about 300 defects to certify the
  conditional bound.
  \item \textbf{E2, score comparison.} S1--S7 under CRC with the same
  splits and seeds. Metric: HRF at fixed $\varepsilon$, plus AUC on the
  pool.
  \item \textbf{Label noise.} Calibrate on labels with a fraction $h$ of
  defects hidden (randomly, or worst-case, i.e.\ the most clean-looking
  first) and measure the realised miss rate against the clean labels. On
  CBRI, calibrate on bootstrap versus gold labels.
\end{itemize}

\subsection{Baselines}

\begin{description}
  \item[B1] Raw softmax: classes in descending probability until their
  mass reaches $1-\alpha$.
  \item[B2] B1 after temperature scaling fitted by NLL.
  \item[B3] Selective classification: answer the top class above the
  most permissive max-probability threshold whose \emph{empirical}
  miscoverage on the tuning rows is $\le\alpha$. In effect this is split
  CP without the $+1$ finite-sample correction.
\end{description}
B2 and B3 fit on the same rows conformal methods calibrate on, so every
method has the same data budget.

\subsection{Metrics}

\begin{itemize}
  \item \textbf{Coverage:} the share of test images whose prediction set
  contains the true label (E1).
  \item \textbf{Realised risk:} the empirical missed-defect rate at
  $\hat\lambda$ on test rows. It must sit at or below $\varepsilon$ on
  average for CRC, and with probability $1-\delta$ for LTT.
  \item \textbf{Exceedance:} the share of seeds whose realised risk
  exceeds $\varepsilon$, i.e.\ what a single deployment experiences.
  \item \textbf{HRF:} the share of test images not auto-cleared. This is
  the cost of the guarantee and the paper's central quantity.
\end{itemize}

\subsection{Compute}

All models run frozen on a single 8\,GB laptop GPU (RTX 5050). Qwen2-VL-2B
runs in bf16 and OWLv2 in fp16, with its preprocessing moved onto the GPU
(the stock CPU resize left the GPU idle about 90\% of the time; outputs
match the stock pipeline to $0.006$). Scoring 3000 images takes 5--25
minutes per score. All conformal computation runs on the CPU in seconds.
